% Chapter 11: The Erd\H{o}s Distinct Distances Problem
\let\LocalMacro\relax

\chapter{The Erd\H{o}s Distinct Distances Problem}

\section{The Problem}

\subsection{Informal}
Scatter a bunch of dots on a piece of paper and look at all the distances between every pair of dots. How many different distances are you guaranteed to see, no matter how you arrange the dots? In 1946, Paul Erd\H{o}s showed that a neat square grid arrangement produces relatively few distinct distances and conjectured that no clever arrangement could do much better. The question stayed open for 64 years until Larry Guth and Nets Katz found a nearly optimal answer using an unexpectedly elegant algebraic trick.

\subsection{Formal}
\begin{equation}
f(n) \geq \frac{n}{\sqrt{\log n}}
\end{equation}

for some constant $c$. This problem sat open for 64 years.

\begin{historicalbox}
Erd\H{o}s was one of the most prolific mathematicians in history, publishing over 1,500 papers. He offered prizes for open problems and could pose beautiful questions on any topic. The distinct distances problem was one of his favorites.
\end{historicalbox}

\section{Mathematical Theory}


\subsection{Prerequisites}
This chapter assumes familiarity with:
\begin{itemize}
\item Basic combinatorics and counting arguments
\item Elementary geometry (distances in the plane)
\item Polynomial algebra (degree, roots)
\end{itemize}

\subsection{The Grid Construction}

\begin{definition}[Integer Grid]
Place $n = m^2$ points at integer coordinates $(i,j)$ for $1 \leq i,j \leq m$. The squared distances are $a^2 + b^2$ for $0 \leq a,b \leq m-1$. The number of distinct squared distances is at most $2m^2 = 2n$, but many of these are not achievable as sums of two squares. A careful analysis shows the number of distinct distances is approximately $n/\sqrt{\log n}$.
\end{definition}

This construction shows $f(n) \leq O(n/\sqrt{\log n})$. The lower bound was the hard part.

\subsection{Previous Bounds}

\begin{itemize}
\item Erd\H{o}s (1946): $f(n) \geq cn^{1/4}$ (using the grid construction and counting) \cite{erdos1946}
\item Chung (1984): $f(n) \geq cn^{2/5}$ \cite{chung1984}
\item Clarkson, Edelsbrunner, Guibas, Stolfi, and Welzl (1990): $f(n) \geq cn^{2/3}/\log n$ \cite{clarkson1990}
\item Solymosi and T\'oth (2003): $f(n) \geq cn^{2/3}/(\log n)^{c}$ \cite{solymosi2003}
\end{itemize}

Progress was incremental and frustratingly slow.

\subsection{The Polynomial Method}

The key new idea came from Dvir's proof of the Kakeya conjecture over finite fields (2008), which used polynomials to count geometric configurations \cite{dvir2009}. The insight: if you have many points with few distances, you can construct a nonzero polynomial that vanishes at all the points, and then use algebraic properties to derive a contradiction.

\begin{definition}[Polynomial Method]
Given a set of points $P \subset \mathbb{R}^2$, find a nonzero polynomial $p(x,y)$ of minimum degree that vanishes at all points in $P$. The degree of $p$ constrains the geometry of $P$.
\end{definition}

\section{The Solution}

\subsection{Guth and Katz (2010)}

Larry Guth and Nets Katz proved in 2010 \cite{guth2015}:

\begin{theorem}[Guth--Katz, 2010 \cite{guth2015}]
For any set of $n$ points in the plane, the number of distinct distances is at least $cn/\log n$ for some constant $c > 0$.
\end{theorem}

This was essentially optimal (matching Erd\H{o}s's grid construction up to the $\log n$ factor).

Their proof combined several ideas:

\begin{enumerate}
\item \textbf{Polynomial partitioning}: Divide the plane into regions using a polynomial surface, then study how points and distances distribute across regions.
\item \textbf{Incidence geometry}: Count how many times points lie on lines (or curves) using the Szemer\'edi--Trotter theorem.
\item \textbf{Joints argument}: Count the number of ``joints'' (points where three non-coplanar lines meet) to control the geometry.
\end{enumerate}

\begin{highlightbox}
The Guth--Katz proof was a masterpiece of combining techniques from different areas. The polynomial method, originally from finite fields, was adapted to the real plane using a novel partitioning argument.
\end{highlightbox}

\subsection{Subsequent Developments}

\begin{itemize}
\item \textbf{Improved bounds}: Later work by Dong, Guth, and others improved the bound slightly.
\item \textbf{Higher dimensions}: The polynomial method has been extended to distinct distances in $\mathbb{R}^d$ for $d \geq 3$.
\item \textbf{Other metric spaces}: The polynomial method has been applied to distinct distances in various settings.
\end{itemize}

\section{Impact}

\begin{itemize}
\item \textbf{Polynomial method}: This proof cemented the polynomial method as a major tool in combinatorial geometry.
\item \textbf{Incidence geometry}: The Guth--Katz approach has influenced many problems in incidence geometry.
\item \textbf{Interdisciplinary connections}: The techniques connect algebra, geometry, and combinatorics in new ways.
\end{itemize}

\section{Exercises}

\begin{exercise}[Basic]
Show that $n$ collinear points determine at most $n-1$ distinct distances. Why doesn't this contradict the theorem?
\end{exercise}

\begin{exercise}[Intermediate]
Construct a set of $n$ points on a circle and count the number of distinct distances. How does this compare to the grid construction?
\end{exercise}

\begin{exercise}[Challenge]
Explain why the polynomial method works better in finite fields than in $\mathbb{R}$. What new difficulties arise in the real case?
\end{exercise}


\begin{exercise}[Computational]
Write a Python/SageMath script to explore a concept from this chapter. See the appendix for starter code.
\end{exercise}

\section{Timeline}

\begin{tabularx}{\linewidth}{lX}
\toprule
\textbf{Year} & \textbf{Event} \\ \midrule
1946 & Erd\H{o}s poses the distinct distances problem \cite{erdos1946} \\ 
1984 & Chung improves the lower bound to $n^{2/5}$ \cite{chung1984} \\ 
2008 & Dvir proves the finite field Kakeya conjecture using polynomials \cite{dvir2009} \\ 
2010 & Guth and Katz prove $f(n) \geq cn/\log n$ \cite{guth2015} \\ 
2015+ & Further improvements and extensions \\ \bottomrule
\end{tabularx}

\section{Citations}

\begin{enumerate}
\item Guth, L., \& Katz, N. (2015). ``On the Erd\H{o}s distinct distances problem in the plane.'' Annals of Mathematics, 181(1), 155--190.
\item Erd\H{o}s, P. (1946). ``On sets of distances of $n$ points.'' American Mathematical Monthly, 53(5), 248--250.
\item Dvir, Z. (2009). ``On the size of Kakeya sets in finite fields.'' Journal of the AMS, 23(4), 1107--1119.
\item Chung, F. (1984). ``The number of different distances determined by $n$ points in the plane.'' Journal of Combinatorial Theory, Series A, 36(3), 342--354.
\item Clarkson, K. L., Edelsbrunner, H., Guibas, L. J., Stolfi, M., \& Welzl, E. (1990). ``Combinatorial complexity bounds for arrangements of curves and spheres.'' Discrete \& Computational Geometry, 5(2), 99--160.
\item Solymosi, J., \& T\'oth, C. D. (2003). ``Distinct distances in the plane.'' Discrete \& Computational Geometry, 30(3), 433--448.
\end{enumerate}
